\scanpage{210}
essayons de déterminer directement les $\varphi_\ast$, endom de la $k$-alg. diff. $\frac12$ simpliciale $A_\ast(k)$.

1) Premier invariant $\varphi_0\in\operatorname{End}_{k\text{-alg}}(S_k)$
\margin{$\begin{gathered}\checkmark\\ -\\ \cdot\\ -\\ \cdot\end{gathered}$}
\deleted{$\varphi_0$} $\varphi_0$ est donné par les
\[
\varphi_0(T^i)=\text{\textcircled{$\tau_i$}}\in S_k\qquad i\geq1
\]
satisfaisant les conditions
\[
\boxed{\tau_i\tau_j=c_{ij}\tau_{i+j}}
\]
Chaque $\varphi_n$ est $\varphi_0$-semi-linéaire.

2) Deuxième invariant $\varphi_1$, endom $\varphi_0$-semi-linéaire de $S_k\{X_0\}[dX_0]$, considéré comme $S_k$-alg. différentielle. \inserted{Il doit être compatible avec les augm. $\alpha$ et $\beta$}
Plus généralement une prop. autonome :

\uline{Prop} Soit $S$, $S'$ deux anneaux comm., $\varphi_0:S\longrightarrow S'$ un hom, on s'intéresse à la détermination des hom d'anneaux différentiels
\[
\varphi_1:\underbrace{S\{X_0\}[dX_0]}_{A_S}\longrightarrow\underbrace{S'\{X_0\}[dX_{\unsure{0|n}}]}_{A_{S'}}
\]
induisant $\varphi_0$ sur $S$, et compatible avec augmentation ($X_0\longmapsto0$, $dX_0\longmapsto0$). Ils sont donnés \unsure{exactmt|exactement} par des systèmes d'éléments
\[
\begin{aligned}
\underset{\substack{\Vert\\A_i+B_i\,dX_0}}{\xi_i}&\in A_{S'}^+=\operatorname{Ker}(\alpha_{S'}:A_{S'}\longrightarrow S')\\
&A_i\in S'\{X_0\}^+,\quad B_i\in S\{X_0\}
\end{aligned}
\]
satisfaisant à $\qquad A_i=A_1^{(i)}$
$\varphi_1$ étant déterminé par les $\xi_i$ via :
\[
\begin{aligned}
\varphi_1(X_0^{(i)})&=\xi_i\\
\varphi_1(dX_0)&=d\xi_1
\end{aligned}
\]
$\varphi$ est compatible avec les graduations (resp. les p.d.) ssi $B_i=0\quad\forall i\geq1$ (resp. $B_i=A_1^{(i-1)}B_1\quad\forall i\geq2$).
% archive sheet number: 209
